\section{Related Work}\label{sec:related}

Following the dual nature of the problem, this section reviews existing methods from simple to complex, covering both domain knowledge on collection operations and AI techniques for agents and selective decisions, and then summarizes them in a SWOT comparison.

% M. Methods
\subsection{Existing Methods}
\textbf{Collection schedules and operations.} Operations research on municipal solid waste collection has focused on the operator's side: vehicle routing, districting, and scheduling~\cite{belien2014}. Published schedules are its resident-facing product. They are easy to understand, but actual arrivals drift with traffic, weather, and variable dwell times along the route. Open fleet feeds such as Taichung's~\cite{taichung_fleet_api} add live positions but leave the interpretation to residents.

\textbf{Arrival-time prediction.} For fixed-route vehicles, arrival times have been predicted with historical averages, regression, Kalman filtering, and neural networks~\cite{chien2002,jeong2004}. These methods generally assume dense location updates and evaluate point accuracy. With sparse observations, simple baselines may be as robust as complex models, and a fair comparison requires time-ordered evaluation that prevents information from the future leaking into training~\cite{tashman2000,hyndman2021}.

\textbf{Pedestrian accessibility.} Walking time depends on the pedestrian network rather than straight-line distance, and walking speed differs substantially across individuals and age groups~\cite{knoblauch1996}. Spatial databases such as PostGIS support indexed nearest-neighbor search and network distance computation for many candidate stops per query~\cite{postgis}.

\textbf{Tool-using LLM agents.} LLMs can be extended with external tools by interleaving reasoning and actions~\cite{yao2023react}, by learning when to call application programming interfaces~\cite{schick2023toolformer}, or by fine-tuning on large API collections~\cite{patil2023gorilla}. Tool use reduces but does not eliminate hallucination: models may still invent arguments, misreport tool outputs, or answer outside their competence~\cite{ji2023hallucination}.

% N. New proposed
\subsection{New Trends and the Proposed Direction}
Two recent trends motivate this work. First, deployed agents are moving from open-ended tool use toward restricted, schema-validated tool sets, trading flexibility for verifiability. Second, selective prediction formalizes when a model should abstain through a risk--coverage trade-off~\cite{chow1970,geifman2017}. GarbageGo combines both: it restricts the LLM to explaining the output of validated domain tools, and it treats departure advice as a selective decision whose primary measure is the false-feasibility rate at a given coverage rather than arrival-time accuracy alone.

% O. Organize
\subsection{Comparative Summary}
Table~\ref{tab:swot} summarizes the strengths, weaknesses, opportunities, and threats of the four method families with respect to departure decisions.

\begin{table*}[t]
\caption{SWOT Comparison of Existing Approaches for Departure Decisions in Curbside Collection}
\label{tab:swot}
\centering
\footnotesize
\begin{tabular}{@{}p{0.15\textwidth}p{0.19\textwidth}p{0.19\textwidth}p{0.19\textwidth}p{0.19\textwidth}@{}}
\toprule
Approach & Strengths & Weaknesses & Opportunities & Threats \\
\midrule
Static schedules and tracking displays~\cite{belien2014,taichung_fleet_api} & Simple, transparent, available citywide & No walking time; no estimate of arrival at the user's stop & Open fleet feeds enable dynamic services & Drift from schedule misleads residents \\
\addlinespace
Arrival-time prediction~\cite{chien2002,jeong2004} & Quantitative ETA with measurable error & Assumes dense data; gives an estimate, not a departure decision & Historical feeds allow backtesting & Ten-minute updates cap accuracy; overfitting \\
\addlinespace
Pedestrian routing~\cite{knoblauch1996,postgis} & Network-based walking time, per-user speed & Ignores truck route and dwell behavior & Spatial databases scale to many stops & Positioning error; age-related speed variation \\
\addlinespace
Tool-using LLM agents~\cite{yao2023react,schick2023toolformer,patil2023gorilla} & Natural-language interaction and flexible orchestration & Hallucinated numbers; unbounded tool use~\cite{ji2023hallucination} & Schema-validated tools and grounding checks & Unsafe advice; out-of-scope answers \\
\bottomrule
\end{tabular}
\end{table*}
